% ---
% id: protocol-compiler/rem:vacuous
% kind: rem
% title: "Party computability read against the readouts"
% summary: >-
%   Read against the joint outcome pair, party computability would hold for
%   every raw-key map; read against the holder's readout it is a genuine
%   restriction. It matters most where both parties' data share one outcome
%   symbol, as when no source is trusted.
% prominence: standard
% status: accepted
% uses: [def:target]
% mentions: [def:validator, transmitter/rem:da1]
% assumes: []
% refs: []
% provenance: newly-derived
% lean:
% ---

\begin{remark}
\label{rem:vacuous}
Read against the joint pair $(x,v)$ instead of against the readouts,
\eqref{eq:partycomp} would be satisfied by \emph{every} $r$ whatsoever, since
$r$ is by definition a function of $(x,v)$; the condition would be vacuous
exactly for the declaration shapes that need it most --- those which place both
parties' data in one outcome symbol, as the $d_A=1$ embedding of
\xref{transmitter/rem:da1} does. Checked against $\rd{\hold}$, it is a genuine
restriction: it fails as soon as two pairs with the same holder readout and the
same announcements carry different key values (check V8 of
Definition~\ref{def:validator}). That embedding --- Alice's measured
components folded into the receiver-side outcome, the \tref{party partition}
naming which components she holds --- is the shape in which a protocol with no
trusted source is declared, and the argument above is why party computability
is checked against the partition in that shape.
\end{remark}
